为确定 $D$ 的符号，先说明其零集。
任取 $x=(s,i)\in\mathcal E$，记其安全端点的易感坐标为 $z=\bar s(x)$。
由命题\ref{prop:reachable}，$h<z<s$ 且 $i=I_z(s)$，所以
\[
D(x)=0
\iff \Theta(s,I_z(s))=\Theta(z,a(z)).
\]
平凡根 $s=z$ 位于安全边界，不属于这里的不安全状态。
由引理\ref{lem:unimodal}，上述等式成立当且仅当
$h<z<e$ 且 $s=\sigma(z)$。
结合已经证明的切换点可达性以及 $i\le K$，得到
\[
\{x\in\mathcal E:D(x)=0\}=\Gamma_K.
\]

在非平凡切换点 $y=(\sigma(z),j(z))$，
沿所属的反向完全跟踪特征线有
\[
\left.\frac{\mathrm d}{\mathrm d s}D(s,I_z(s))\right|_{s=\sigma(z)}
=-\left.\frac{\mathrm d}{\mathrm d s}\Theta(s,I_z(s))\right|_{s=\sigma(z)}>0,
\]
其中严格符号来自非平凡根位于单峰函数的严格下降段。
由于 $q=1$ 时 $\dot s=-csi<0$，故 $L_{f_1}D(y)<0$，特别地
$\nabla D(y)\ne0$。
端点映射的光滑性给出 $D$ 在各有关点附近的光滑性；
在容量切换点，$a(z)>0$ 及 $G'(z)\ne0$ 还给出环境开邻域上的光滑延拓，
这里只将其用于求导。

另一方面，$\chi(s,i)=i-I_\Gamma(s)$ 满足
$\nabla\chi=(-I_\Gamma'(s),1)\ne0$。
沿 $\Gamma_K$ 上的恒等式 $D(s,I_\Gamma(s))=0$ 求导，
得到 $\nabla D\cdot(1,I_\Gamma')=0$；
在容量端点通过 $s\uparrow s_B$ 取极限，关系同样成立。
因此存在标量 $\lambda(y)$，使
\[
\nabla D(y)=\lambda(y)\nabla\chi(y),\qquad
\lambda(y)=\frac{L_{f_1}D(y)}{L_{f_1}\chi(y)}<0.
\]
再由引理\ref{lem:strict-crossing}知 $L_{f_0}D(y)<0$。
沿自然轨道的时间方向有 $\dot s<0$，所以在触点 $s=\eta$ 处，
$D$ 沿该轨道关于 $s$ 的导数严格为正。
向未来即向较小 $s$ 移动时，$D$ 因而先变为负值。
从 $y_\Gamma$ 到 $x$ 的自然轨道段已知保持在 $\mathcal E$ 内，
而该轨道没有第二个非平凡切换交点。
由零集刻画和连续性，到 $x$ 为止始终有 $D<0$。
